\section{Related Work}

\subsection{Repository-Level Code Generation}

Recent LLM-based software engineering systems increasingly target repository-level generation.
ChatDev~\cite{qian2023chatdev} and MetaGPT~\cite{hong2023metagpt} decompose development into role-based multi-agent workflows, while SWE-agent~\cite{yang2024sweagent}, OpenHands~\cite{wang2024openhands}, and CodeAgent~\cite{zhang2024codeagent} equip agents with tools for navigating, editing, and testing real codebases.
AutoP2C~\cite{shi2025autop2c} is closest to our setting, transforming academic papers into executable repositories.
Yet these systems optimize for executable completeness: implementation gaps are defects to eliminate.
A2P targets a different objective, where the repository must run while preserving deliberately chosen gaps as learner-owned work.

\subsection{Educational Content Generation}

LLMs have also been used to generate educational materials and learning interactions.
Systems such as PPTAgent~\cite{zheng2025pptagent}, DeepPresenter~\cite{zheng2026deeppresenter}, STORM~\cite{shao2024storm}, and PustakAI~\cite{sharma2025pustakai} produce instructional artifacts such as slides, articles, and textbooks, while LessonPlanner~\cite{wang2024lessonplanner} supports pedagogy-driven lesson planning.
Tutoring systems take a complementary approach: SimClass~\cite{zhang2024simclass} simulates classroom interaction, MedTutor~\cite{jang2026medtutor} supports case-based medical education, and EDF~\cite{cohn2026edf} implements theory-driven adaptive scaffolding.
Co-design work with K-12 educators further shows that project design is a major pain point where LLM support is desired~\cite{lee2025codesigning}.
However, these systems primarily generate instructional content or interactions around learning.
A2P instead generates the executable setting in which learning occurs: source concepts are instantiated as files, tests, tasks, scaffolds, and learner-owned implementation gaps.

\subsection{Scaffolding and Learner Modeling}

Scaffolding theory explains why incompletion in a learning repository must be controlled rather than arbitrary.
Vygotsky's zone of proximal development~\cite{vygotsky1978} and Wood et al.'s scaffolding framework~\cite{wood1976scaffolding} establish that effective instruction targets work learners cannot yet complete independently but can complete with support.
Cognitive load theory~\cite{sweller1988cognitive} provides the complementary constraint: excessive unsupported complexity prevents learning, while productive failure~\cite{kapur2008productive} and desirable difficulties~\cite{bjork2011desirable} show that appropriately challenging struggle can improve learning.
This balance is central to PBL, where authentic artifacts require structured support rather than minimal guidance~\cite{kirschner2006minimal,hmelosilver2007scaffolding,krajcik2014pbl}.
A2P operationalizes this balance as a repository-generation problem: deciding which components should be complete, scaffolded, or left as learner-owned implementation gaps.

On the computational side, learner modeling has evolved from Bayesian Knowledge Tracing~\cite{corbett1995knowledge} through deep knowledge tracing~\cite{piech2015deep} to LLM-augmented approaches such as SINKT~\cite{gao2024sinkt}, which uses LLM-encoded concept semantics for inductive knowledge tracing.
These systems model learners for \emph{runtime} adaptation---selecting or sequencing existing content.
A2P instead uses a one-shot learner profile to condition \emph{generation}, deciding what to scaffold, what to leave as productive struggle~\cite{kapur2008productive}, and what to omit as noise \emph{before} the project exists.
